\subsection{KRULL 整环}

定义 3. 若整环 A 存在一个族

\((v_{\iota})_{\iota\in I}\) 的赋值，它们定义在 A 的分式域 \(\mathbf{K}\) 上并满足以下性质：

(AK \(_{I}\) ) 赋值 \(v_{\iota}\) 是离散的；

(AK \(_{II}\) ) 各个 v \(_{\iota}\) 的赋值环的交为 A；

\((\mathbf{A}\mathbf{K}_{\mathrm{III}})\) 对所有 \(x \in \mathbf{K}^*\)，指标 \(\iota \in \mathbf{I}\) 中满足 \(v_{\iota}(x) \neq 0\) 的集合是有限的。

显然，只需对 A - (0) 中的元素 x 验证条件 \((\mathrm{AK}_{\mathrm{III}})\) 即可。

\textbf{例子}\par

\begin{enumerate}
\def\labelenumi{(\arabic{enumi})}
\item
  每个离散赋值环都是 Krull 整环。
\item
  更一般地，每个主理想整环 A 都是 Krull 整环。事实上，设 \((p_{\iota})_{\iota\in I}\) 是 A 的极端元代表系，令 \(v_{\iota}\) 为由 \(p_{\iota}\) 定义在 A 的分式域上的赋值（第 VI 章，§ 3，第 3 节，例 4）。立即可见，族 \((v_{\iota})_{\iota\in I}\) 满足性质 \((\mathbf{AK}_{\mathbf{I}}), (\mathbf{AK}_{\mathbf{II}})\) 和 \((\mathbf{AK}_{\mathbf{III}})\)。
\item
  设 F 是一个域，\((\mathbf{R}_{i})_{1\leqslant i\leqslant n}\) 是 F 的有限个子环组成的族，且这些子环都是 Krull 整环。则它们的交 \(S=\bigcap_{j=1}^{n}R_{j}\) 是 Krull 整环。对 \(1\leqslant j\leqslant n\)，设 \((v_{j\iota})_{\iota\in I_j}\) 是定义在 \(R_{j}\) 的分式域上、满足 \((\mathrm{AK}_{\mathrm{I}}),(\mathrm{AK}_{\mathrm{II}}),(\mathrm{AK}_{\mathrm{III}})\) 的赋值族（其中以 \(R_{j}\) 代替 A）。记 \(w_{j\iota}\) 为 \(v_{j\iota}\) 在 S 的分式域上的限制。则族 \((v_{j\iota})_{1\leqslant j\leqslant n,\iota\in I_j}\) 显然满足 \((\mathrm{AK}_{\mathrm{II}})\)（其中以 S 代替 A），并且也满足 \((\mathrm{AK}_{\mathrm{III}})\)，因为指标 j 的集合是有限的。赋值 \(w_{j\iota}\) 或者是离散的，或者是不正规的。只保留其中离散的赋值，便得到一个显然满足 \((\mathrm{AK}_{\mathrm{I}}),(\mathrm{AK}_{\mathrm{II}})\) 和 \((\mathrm{AK}_{\mathrm{III}})\) 的族（其中以 S 代替 A）。因此 S 确实是 Krull 整环。
\item
  特别地，如果 A 是 Krull 整环，且 \(K'\) 是 A 的分式域 K 的子域，则 \(K' \cap A\) 是 Krull 整环。
\end{enumerate}

定理 2. 设 A 是整环。A 为 Krull 整环的充要条件是满足以下两个条件：

\begin{enumerate}
\def\labelenumi{(\alph{enumi})}
\item
  A 是完全整闭的；
\item
  \(\mathbf{A}\) 的每个非空除性整理想族都有一个极大元（相对于关系 \(\subset\)）。
\end{enumerate}

此外，若 \(\mathrm{P}(\mathrm{A})\) 是 \(\mathrm{D}(\mathrm{A})\) 的极端元集合，则 \(\mathrm{P}(\mathrm{A})\) 是 Z-模 \(\mathrm{D}(\mathrm{A})\) 的一组基，而 \(\mathrm{D}(\mathrm{A})\) 的正元正是以 \(\mathrm{P}(\mathrm{A})\) 的元素为线性组合、系数满足 \(\geqslant0\) 的那些元素。

设 A 是 Krull 整环。它是完全整闭的（第 VI 章，§ 4，第 5 节，命题 9 的推论）。设 \((v_{\iota})_{\iota\in I}\) 是定义在 A 的分式域 K 上、满足 \((\mathrm{AK}_{\mathrm{I}})\)、\((\mathrm{AK}_{\mathrm{II}})\) 和 \((\mathrm{AK}_{\mathrm{III}})\) 的赋值族。可设 \(v_{\iota}\) 都是赋范的（第 VI 章，§ 3，第 6 节，定义 3）。对于所有 \(\mathfrak{a} \in \mathrm{I}(\mathbf{A})\)，记：

\[
v _ {\iota} (\mathfrak{a}) = \sup _ {\mathfrak{a} \subset Ax} (v _ {\iota} (x));\tag{3}
\]

于是 \(v_{\iota}(\mathfrak{a}) \in \mathbf{Z}\)。事实上，若 \(a\) 是 \(\mathfrak{a}\) 的非零元，关系 \(Ax \supset Aa\) 蕴含 \(v_{\iota}(x) \leqslant v_{\iota}(a)\)（由 \((\mathrm{AK}_{\mathrm{II}})\)），这说明 \(v_{\iota}(x)\)（\(\mathfrak{a} \subset Ax\)）组成的族有上界。我们建立以下性质：

\begin{enumerate}
\def\labelenumi{(\arabic{enumi})}
\tightlist
\item
  设 \(\mathfrak{a}\) 是除性分式理想；\(y \in \mathfrak{a}\) 的充要条件是 \(v_{\iota}(y) \geqslant v_{\iota}(\mathfrak{a})\) 对所有 \(\iota \in \mathbf{I}\) 都成立。
\end{enumerate}

由于 \(\mathfrak{a}\) 是除性的，关系 \(y \in \mathfrak{a}\) 等价于关系“\(\mathfrak{a} \subset \mathbf{A}x\) 蕴含 \(y \in \mathbf{A}x\)”。现在，由 \((\mathrm{AK}_{\mathrm{II}})\)，关系 \(y \in \mathbf{A}x\) 等价于“\(v_{\iota}(y) \geqslant v_{\iota}(x)\) 对所有 \(\iota \in \mathbf{I}\) 都成立”。从而得到 (1)。

\begin{enumerate}
\def\labelenumi{(\arabic{enumi})}
\setcounter{enumi}{1}
\tightlist
\item
  设 \(\mathfrak{a}\) 和 \(\mathfrak{b}\) 是 \(\mathbf{A}\) 的两个除性分式理想；\(\mathfrak{a} \subset \mathfrak{b}\) 的充要条件是 \(v_{\iota}(\mathfrak{a}) \geqslant v_{\iota}(\mathfrak{b})\) 对所有 \(\iota \in \mathbf{I}\) 都成立。
\end{enumerate}

这立即由性质 (1) 得出。

\begin{enumerate}
\def\labelenumi{(\arabic{enumi})}
\setcounter{enumi}{2}
\tightlist
\item
  若 \(x \in \mathbf{K}^*\)，则 \(v_{\iota}(\mathbf{A}x) = v_{\iota}(x)\)。
\end{enumerate}

若 \(A y \supset A x\)，则 \(v_{\iota}(y) \leqslant v_{\iota}(x)\)（由 \((\mathbf{A K}_{\mathrm{II}})\)），且 \(v_{\iota}(y)\) 的最小值在 \(y = x\) 处取得。

\begin{enumerate}
\def\labelenumi{(\arabic{enumi})}
\setcounter{enumi}{3}
\tightlist
\item
  对所有 \(\mathfrak{a} \in \mathbf{I}(\mathbf{A})\)，指标 \(\iota \in \mathbf{I}\) 中满足 \(v_{\iota}(\mathfrak{a}) \neq 0\) 的数目是有限的。
\end{enumerate}

存在 \(x, y\)（属于 \(\mathbf{K}^*\)），使得 \(Ax \subset \mathfrak{a} \subset Ay\)。由性质 (2) 和 (3)，\(v_{\iota}(x) \geqslant v_{\iota}(\mathfrak{a}) \geqslant v_{\iota}(y)\) 对所有 \(\iota \in I\) 都成立。于是只需应用 \((\mathrm{AK}_{\mathrm{III}})\) 即可。

因此我们证明了以下引理：

引理 1. 若 A 是 Krull 整环，且 \((v_{\iota})_{\iota \in I}\) 是定义在 K 上、满足 \((\mathrm{AK}_{\mathrm{I}})\)、\((\mathrm{AK}_{\mathrm{II}})\) 和 \((\mathrm{AK}_{\mathrm{III}})\) 的赋范赋值族，则映射 \(\mathfrak{a} \mapsto (v_{\iota}(\mathfrak{a}))_{\iota \in I}\) 将 A 的除性整理想集合（按 \(\subset\) 排序）递减且单射地映入有序群直和 \(\mathbf{Z}^{(I)}\) 的正元集合。

既然如此，\(\mathbf{Z}^{(I)}\) 的每个非空正元集合都有极小元（《代数》第 VI 章，§ 1，第 13 节，定理 2）。因此 A 确实满足该命题的性质 (b)。

反之，设 A 是满足该命题性质 (a) 和 (b) 的整环。由于 A 完全整闭，D(A) 是有序群（第 2 节，定理 1）。这个群是格（第 1 节，命题 2）。由该命题的条件 (b)，D(A) 的每个非空正元族都有极小元。设 P(A) 是 D(A) 的极端元集合。于是（《代数》第 VI 章，§ 1，第 13 节，定理 2）P(A) 是 Z-模 D(A) 的一组基，而 D(A) 的正元是 P(A) 的元素以正整数为系数的线性组合。
 
因此，对于 \(x \in \mathbf{K}^*\)，可以通过下式定义有理整数 \(v_{\mathbb{P}}(x)\)（其中 \(\mathbf{P} \in \mathbb{P}(\mathbf{A})\)）：

\[
\operatorname{div} (x) = \sum_ {\mathbf {P} \in \mathbf {P} (\mathbf {A})} v _ {\mathbf {P}} (x) \cdot \mathbf {P}.\tag{4}
\]

并记 \(v_{\mathbf{P}}(0) = +\infty\)。

由以下关系

\[
\operatorname{div} (x y) = \operatorname{div} (x) + \operatorname{div} (y)
\]

以及

\[
\operatorname{div} (x + y) \geqslant \inf (\operatorname{div} (x), \operatorname{div} (y)),
\]

对于 \(x, y\) 及 \(x + y\)（均属于 \(\mathbf{K}^*\)），我们推出 \(v_{\mathbb{P}}\) 是 \(\mathbf{K}\) 上的离散赋值。\(x \in \mathbf{A}\) 的充要条件是 \(\mathrm{div}(x) \geqslant 0\)，也就是 \(v_{\mathbb{P}}(x) \geqslant 0\) 对所有 \(\mathbb{P} \in \mathbb{P}(\mathbf{A})\) 都成立。因此，\(v_{\mathbb{P}}\) 满足条件 \((\mathbf{AK}_{\mathbf{I}})\) 和 \((\mathbf{AK}_{\mathbf{II}})\)，显然也满足 \((\mathbf{AK}_{\mathbf{III}})\)。

推论。诺特环为 Krull 整环的充要条件是它是整闭整环。

整闭的诺特整环是完全整闭的（第 V 章，§ 1，第 4 节）。

存在非诺特的 Krull 整环，例如域 K 上、具有无穷多个不定元的多项式环 \(K[X_{n}]_{n\in\mathbf{N}}\)（参见习题 8）。

